\section{Conclusion}
\label{sec:conclusion}

The economic consequences of advanced AI depend on more than whether machines can
perform existing jobs. A distinct transition occurs when AI can perform the research
required to produce better AI. The toy model represents this transition as a gradual
change in the composition of research. When human and automated-research services are
gross substitutes, the automated share rises with the research wage relative to its
unit resource cost.

Population growth matters while human research dominates. With $\phi<1$ and a
constant research-employment share, capability growth is semi-endogenous and equals
$\eta n/(1-\phi)$. When automated research dominates, capability growth is tied
instead to the growth of automated-research services. In the Cobb--Douglas production
benchmark, a constant-growth path exists if
$1-\phi-\eta\beta/(1-\alpha-\beta)>0$. If this condition fails, the feedback between
output and automated research is too strong for a positive constant-growth solution.

The production block also separates the wage response from the response of the
production-labor share. An expansion of AI services reduces this share when
$\sigma>1$, but it raises the wage when $\sigma<1/\alpha$. Hence both outcomes occur
when $1<\sigma<1/\alpha$. This conditional result is a benchmark, not the paper's
standalone contribution. In general equilibrium, capital deepening and the
reallocation of workers between production and research also affect wages and the
aggregate labor share.

The production elasticity determines whether the economy reaches that asymptotic
research regime. When $\sigma<1$, labor remains essential and the productive return
to further reductions in the unit cost of AI services vanishes. When $\sigma=1$,
positive wage growth makes automated research asymptotically dominant. When
$\sigma>1$, the production-labor share converges to zero and the return to capital
rises with capability, so a conventional balanced-growth path with positive
capability growth cannot maintain a constant capital--output ratio.

Market power creates a genuine policy trade-off. The monopolist restricts the current
supply of AI services, but its operating rents raise the private return to frontier
research. Because it does not appropriate all consumer surplus and knowledge
benefits, private research remains insufficient whenever $Q_t>q_t$. Regulating the
AI-service price without replacing the lost innovation return can therefore improve
current deployment while slowing capability growth. The natural policy comparison is
between laissez-faire monopoly, regulated access, explicit research support, and a
joint policy that targets both margins.
